\usepackage{listings} % библиотека листингов
\usepackage{color} % подсветка листинга

\definecolor{mygreen}{rgb}{0,0.6,0}
\definecolor{mymauve}{rgb}{0.58,0,0.82}

% \usepackage{minted} % не используется в работе (везде lstlisting), требует sudo для установки в TL2026
\usepackage{bold-extra}

\lstset{
          basicstyle=\footnotesize\ttfamily,
          aboveskip=11pt,
          belowskip=11pt,
          language=python,
          keywordstyle=\color{blue},
          deletekeywords={local},
          breaklines=true,
          commentstyle=\color{teal},
          breakatwhitespace=false,
          showspaces=false,
          showstringspaces=false,
          escapeinside={(*}{*)},
          % mathescape=false: оставляем минусы в коде как ASCII-дефис,
          % не превращая в матрежим (иначе - подменяется на Unicode-минус).
          mathescape=false,
          extendedchars=true,
          % literate: явная подмена дефиса на дефис толщиной 1 — заставляет
          % listings рендерить «-» как ASCII-символ, что предотвращает
          % замену на Unicode-минус.
          literate={-}{{-}}1
    }
\usepackage{lstfiracode}

% Подписи к листингам на русском языке.
\renewcommand\lstlistingname{Листинг}
\renewcommand\lstlistlistingname{Листинги}



